\section{Experiments}

\begin{table*}[!t]
\centering
% \small
\setlength{\tabcolsep}{4pt}
\renewcommand{\arraystretch}{1.08}
\resizebox{\textwidth}{!}{
\begin{tabular}{
lccccc
@{\hspace{0.6em}} c @{\hspace{0.6em}}
lccccc
}
\toprule

\multicolumn{6}{>{\columncolor{lightpink}}c}{\textbf{Qwen3-8B}}
&
&
\multicolumn{6}{>{\columncolor{lightpink}}c}{\textbf{Mistral-7B-v0.3}} \\
\cmidrule(lr){1-6} \cmidrule(lr){8-13}

\textbf{Method}
& \textbf{ALCE} & \textbf{ScholarQA} & \textbf{QASPER} & \textbf{LAiW} & \textbf{LegalBench}
&
&
\textbf{Method}
& \textbf{ALCE} & \textbf{ScholarQA} & \textbf{QASPER} & \textbf{LAiW} & \textbf{LegalBench} \\

\midrule
Initial & 83.33 & 66.36 & 53.25 & 34.14 & 74.93
&
&
Initial & 68.58 & 27.54 & 54.45 & 20.82 & 53.3 \\

ForwardGen & 85.18 & 65.42 & 51.84 & 33.52 & 74.95
&
&
ForwardGen & 62.41 & 48.65 & 51.34 & 32.15 & 60.27 \\

Artifact-Only & 87.86 & 67.18 & 52.52 & 34.69 & 74.93
&
&
Artifact-Only & 63.64 & 55.83 & 50.21 & 34.82 & 62.22 \\

WebGLM & 84.52 & 65.84 & 53.36 & 36.61 & \textbf{76.12}
&
&
WebGLM & 65.11 & 53.27 & 50.81 & 39.41 & 61.25 \\

WebCPM & 46.73 & 60.27 & 49.89 & 23.23 & 74.92
&
&
WebCPM & 57.90 & 50.41 & 54.40 & 38.68 & 63.10 \\

\textbf{RetroGen (Ours)} & \textbf{88.03} & \textbf{67.71} & \textbf{56.15} & \textbf{37.65} & 75.37
&
&
\textbf{RetroGen (Ours)} & \textbf{68.93} & \textbf{61.12} & \textbf{56.27} & \textbf{41.29} & \textbf{65.98} \\

\midrule

\multicolumn{6}{>{\columncolor{lightpink}}c}{\textbf{Qwen2.5-7B}}
&
&
\multicolumn{6}{>{\columncolor{lightpink}}c}{\textbf{Olmo-3-1025-7B}} \\
\cmidrule(lr){1-6} \cmidrule(lr){8-13}

\textbf{Method}
& \textbf{ALCE} & \textbf{ScholarQA} & \textbf{QASPER} & \textbf{LAiW} & \textbf{LegalBench}
&
&
\textbf{Method}
& \textbf{ALCE} & \textbf{ScholarQA} & \textbf{QASPER} & \textbf{LAiW} & \textbf{LegalBench} \\

\midrule

Initial & 85.37 & 56.42 & 47.03 & 34.71 & 72.16
&
&
Initial & 90.70 & 53.86 & 55.17 & 39.24 & 63.57 \\

ForwardGen & 86.27 & 58.36 & 49.62 & 34.18 & 71.32
&
&
ForwardGen & 89.34 & 55.42 & 53.62 & 37.42 & 64.84 \\

Artifact-Only & 87.48 & 62.74 & 50.04 & 33.39 & 72.41
&
&
Artifact-Only & 91.42& 58.17 & 54.36 & 38.27 & 66.18 \\

WebGLM & 84.26 & 59.85 & 50.15 & 35.24 & 71.27 
&
&
WebGLM & 91.27 & 56.78 & 52.97 & 36.61 & 64.25\\

WebCPM & 67.75 & 54.27 & 48.85 & 31.95 & 72.80
&
&
WebCPM & 74.17 & 53.41 & 51.84 & 34.71 & 64.25 \\

\textbf{RetroGen (Ours)} & \textbf{90.50} & \textbf{64.18} & \textbf{51.64} & \textbf{36.00} & \textbf{75.01}
&
&
\textbf{RetroGen (Ours)} & \textbf{92.83} & \textbf{61.78} & \textbf{56.42} & \textbf{39.65} & \textbf{67.30} \\

\bottomrule
\end{tabular}
}
\caption{Main results across four open-source backbone models. 
\textsc{RetroGen} achieves the best average performance across evidence-grounded generation and domain-specific reasoning benchmarks.}
% \vspace{-4.5mm}
\label{tab:main_results}
\end{table*}
\begin{table*}[!t]
\centering
\small
\setlength{\tabcolsep}{6pt}
\renewcommand{\arraystretch}{1.05}
\begin{tabular}{lcccccc}
\toprule
\textbf{Filter} & \textbf{ALCE} & \textbf{ScholarQA} & \textbf{QASPER} & \textbf{LAiW} & \textbf{LegalBench} & \textbf{Avg} \\
\midrule
Initial & 83.33 & 66.36 & 53.25 & 34.14 & 74.93 & 62.40 \\
only $s_{\mathrm{qual}}$ & 84.88 & 66.62 & 53.48 & 34.72 & 75.46 & 63.03 \\
only $s_{\mathrm{con}}$ & 85.42 & 66.18 & 53.86 & 35.64 & 75.22 & 63.26 \\
next-$N$ replace & 86.42 & 66.74 & 54.18 & 35.86 & 75.12 & 63.66 \\
only $s_{\mathrm{grd}}$ & 86.18 & 66.52 & 55.42 & 36.48 & 75.44 & 64.01 \\
only $s_{\mathrm{rub}}$ & 88.18 & 67.82 & 55.28 & 36.05 & 74.96 & 64.46 \\
\textbf{RetroGen} ($\lambda{=}0.25$) & 88.03 & 67.71 & \textbf{56.15} & \textbf{37.65} & 75.37 & \textbf{64.98} \\
\bottomrule
\end{tabular}
\caption{Verification-dimension ablation on Qwen3-8B under a matched token budget. Uniformly weighted scoring yields the best average; single-signal filters and next-$N$ replacement remain above the initial checkpoint but are less stable.}
\label{tab:verifier_ablation}
\end{table*}


\subsection{Experimental Setup}
\label{sec:setup}

\paragraph{Domains and Artifacts.}
We evaluate \textsc{RetroGen} on three evidence-grounded long-form generation domains: 
(i) \textbf{scientific writing}, where the task is to draft a related-work section conditioned on a paper abstract and the expert artifact $y^*$ is the author-written related-work section from the corresponding paper;
(ii) \textbf{financial analysis}, where the agent produces an analyst-style report grounded in U.S. SEC 10-K filings from the EDGAR corpus~\citep{loukas2021edgar} and $y^*$ is a human-written financial narrative; and
(iii) \textbf{legal judgment drafting}, where the agent drafts a first-instance judgment grounded in statutes and prior cases, with $y^*$ given by the official adjudicated judgment.
These domains stress-test RetroGen under distinct evidence regimes: citation-heavy scientific synthesis, document-grounded financial reasoning, and statute/case-grounded legal argumentation.

\paragraph{Tools and Execution Environment.}
All trajectories are generated using domain-specific instantiations of the abstract action space $\mathcal{A}$ defined in Section~\ref{sec:method}. To ensure that supervision reflects executable evidence-seeking behavior, observations are populated through actual tool calls rather than model-imagined snippets. We cap each trajectory at $30$ tool steps and maintain an $8{,}000$-character sliding context window.

\paragraph{Backbone Models.}
We evaluate \textsc{RetroGen} on four open-source LLMs: Qwen3-8B~\citep{yang2025qwen3}, Qwen2.5-7B~\citep{Yang2024Qwen25TR}, Mistral-7B-v0.3~\citep{Jiang2023Mistral7}, and Olmo-3-1025-7B~\citep{olmo2025olmo}. 

\paragraph{Training Corpus.}
For each backbone, \textsc{RetroGen} reconstructs candidate trajectories from expert artifacts and filters them using the multi-dimensional verification score in Section~\ref{sec:method} with uniform weights $\lambda=0.25$. Thresholds are calibrated per domain from the score distribution (0.52 / 0.48 / 0.63 for scientific / financial / legal; Appendix~\ref{app:verification}), retaining approximately $25$K verified trajectories. We construct a $\sim$50M-token SFT corpus with a fixed token-budgeted mixture: $80\%$ verified agentic trajectories, split $40$:$20$:$20$ across scientific, financial, and legal domains, and $20\%$ general instruction-following, math, coding, and multi-turn dialogue data. This mixture preserves general capabilities while emphasizing evidence-seeking behavior.

\paragraph{Baselines and Ablations.}
We compare against controlled alternatives under the same token budget, sequence length, and optimization configuration:
\begin{enumerate}[label=(\arabic*), leftmargin=*, topsep=3pt, partopsep=0pt, itemsep=0pt, parsep=3pt]
    \item \textbf{Initial}: the checkpoint before training.
    \item \textbf{ForwardGen}: trajectories generated by the same backbone from the task specification $\hat{x}$ without access to the expert artifact, isolating the value of target-anchored retrospective reconstruction.
    \item \textbf{Artifact-Only}: the same reconstructed dataset with all intermediate tool interactions removed, isolating trajectory-level process supervision from input-output supervision.
    \item \textbf{Public-Baseline}: the $80\%$ agentic slice is replaced with WebGLM-QA~\citep{liu2023webglm} and WebCPM-WK~\citep{qin2023webcpm}, two public citation-grounded long-form generation corpora that provide retrieve-then-write supervision distilled from stronger systems.
\end{enumerate}

\paragraph{Evaluation.}
We evaluate along three axes. First, \emph{evidence-grounded long-form generation} is measured by ALCE~\citep{gao2023enabling}, ScholarQA~\citep{openscholar}, and QASPER~\citep{mcdonald2022detect}. Second, \emph{domain-specific grounded reasoning} is assessed with legal reasoning benchmarks including LAiW~\citep{dai2025laiw} and LegalBench~\cite{guha2023legalbench}. Third, \emph{general capability retention} is measured by GSM8K~\citep{cobbe2021gsm8k}, MMLU-Pro~\citep{wang2024mmlu}, IFEval~\citep{zhou2023instruction}, and HumanEval~\citep{chen2021codex}. All evaluations use vLLM-backed zero/few-shot inference~\citep{kwon2023efficient}, with details deferred to Appendix.

\subsection{Main Results}
\begin{figure*}[!t]
    \centering
    \includegraphics[width=0.8\linewidth]{figures/exp_astabench_overview.pdf}
    \caption{AstaBench literature-understanding performance on LitQA2 and SQA. }
    \label{fig:exp_asta}
\end{figure*}

\begin{figure*}[!t]
    \centering
    \includegraphics[width=0.8\linewidth]{figures/exp_transitions.pdf}
    \caption{AstaBench action-transition matrices, where each cell denotes the probability of transitioning from the current action to the next action.}
    \label{fig:exp_trans}
    \vspace{-3mm}
\end{figure*}

Table~\ref{tab:main_results} reports the performance of \textsc{RetroGen} and all baselines across the four backbone models. We summarize three main observations.

\begin{itemize}[leftmargin=*]
    \item \textbf{Retrospective anchoring improves evidence-grounded generation.}
Across ALCE, ScholarQA, and QASPER, \textsc{RetroGen} consistently outperforms the Initial and ForwardGen models across backbones. Compared with ForwardGen, which generates trajectories from the task specification $\hat{x}$ without access to the expert artifact, \textsc{RetroGen} improves the average evidence-grounded generation score by +4.8 points, with gains of +4.3 on ALCE, +6.7 on ScholarQA, and +3.5 on QASPER. This suggests that anchoring trajectory reconstruction on expert artifacts helps reduce unguided exploration and produces more reliable evidence-seeking behavior.

\item \textbf{Trajectory-level supervision provides gains beyond final-artifact imitation.}
The comparison with Artifact-Only isolates the value of intermediate tool-interaction traces. On domain-specific reasoning benchmarks, \textsc{RetroGen} improves over Artifact-Only by +2.7 points on average, with gains of +3.4 on LAiW and +2.0 on LegalBench. This indicates that final artifacts alone do not fully expose the procedural structure needed for grounded reasoning; reconstructing how evidence is searched, selected, and organized provides additional supervision for complex synthesis tasks.

    \item \textbf{Verified agentic trajectories are more effective than collapsed retrieve-then-write supervision.}
    \textsc{RetroGen} also outperforms the Public-Baseline, where the same token budget is allocated to public web-grounded long-form QA data in a retrieve-then-write format. The gains are most pronounced on tasks requiring multi-step evidence aggregation, suggesting that supervision over explicit tool-use trajectories teaches the model more than observing only a retrieved context and final answer.
\end{itemize}

\subsection{In-Depth Analysis}

\subsubsection{Verification-Dimension Ablation}
\label{sec:verifier_ablation}

The main results use a uniformly weighted combination of the four verification signals. To test whether this combination is necessary, we retrain Qwen3-8B under the same token budget while filtering with a single score, or while replacing the selected traces with the next-highest-scoring ones (\emph{next-$N$ replace}).
Table~\ref{tab:verifier_ablation} shows that the weighted verifier achieves the best average. Single-dimension filters remain above the initial checkpoint, so moderately scored traces can still help; however, next-$N$ replace underperforms full \textsc{RetroGen}, indicating that usefulness is graded and that the threshold trades quality against data volume.

\subsubsection{Dynamic Agentic Evaluation on AstaBench}
\label{sec:astabench}
\begin{figure}[!t]
    \centering
    \includegraphics[width=\linewidth]{figures/fig_general_capability.pdf}
    \caption{General capability retention across four backbone models. Dumbbell markers compare each initial checkpoint with its \textsc{RetroGen}-trained counterpart.}
    \label{fig:general_cap}
    \vspace{-3mm}
\end{figure}

Standard evidence-grounded generation benchmarks evaluate the quality of final
answers, but they only partially reveal whether a model can behave as a
reliable research agent. To further examine the process-level behavior learned
by \textsc{RetroGen}, we evaluate on AstaBench~\citep{bragg2025astabench}, a
dynamic agentic benchmark for scientific research tasks. Unlike static QA
benchmarks, AstaBench requires the model to iteratively search for papers,
inspect retrieved content, and decide when sufficient evidence has been
collected. This makes it well suited for testing whether retrospective process
supervision improves the agent's evidence-seeking policy rather than only its
final-output style.

\paragraph{Benchmark and protocol.}
We focus on the literature-understanding tasks LitQA2-Validation and SQA (details in Appendix~\ref{app:eval}).
The initial model and its \textsc{RetroGen}-trained
counterpart use the same backbone, tool inventory, sampling configuration, and
maximum interaction budget (\texttt{max\_rounds}$=15$).

\paragraph{Performance on literature understanding.}
Figure~\ref{fig:exp_asta} shows that \textsc{RetroGen} improves most
AstaBench literature-understanding metrics. On LitQA2, \textsc{RetroGen}
doubles accuracy from $0.20$ to $0.40$ and improves precision from $0.33$ to
$0.67$, while maintaining the same coverage of $0.60$. This indicates that the
model becomes more accurate and selective without simply abstaining more often.
On SQA, \textsc{RetroGen} improves the global average from $0.47$ to $0.59$.
The largest gains come from citation-related metrics: citation precision
increases from $0.47$ to $0.67$, and citation recall increases from $0.33$ to
$0.57$. Answer precision remains unchanged at $0.82$, suggesting that the main
benefit is not merely more fluent answer generation, but better grounding and
attribution of claims to supporting evidence.

\paragraph{Action-transition analysis.}
To understand where these gains come from, we further analyze the action
transitions recorded by the AstaBench sandbox. We collapse raw tool calls into
three high-level action types: \textsc{Search}, \textsc{Read}, and \textsc{END}.
Figure~\ref{fig:exp_trans} compares the empirical transition matrices of the
initial model and \textsc{RetroGen}. The initial model frequently repeats search
actions after searching, with a search-to-search probability of $0.54$,
suggesting redundant exploration. In contrast, \textsc{RetroGen} reduces this
probability to $0.39$ and increases the search-to-read transition from $0.15$
to $0.23$, indicating that it is more likely to inspect retrieved evidence
rather than continue issuing new searches. After reading, \textsc{RetroGen}
also transitions to \textsc{END} much more often, increasing the read-to-end
probability from $0.20$ to $0.55$.

Together, these patterns suggest that \textsc{RetroGen} learns a more
goal-directed research policy: it searches to identify candidate evidence,
reads to verify and extract support, and terminates once sufficient evidence has
been gathered. This provides process-level evidence for our main claim that
retrospective supervision improves the structure of evidence-seeking behavior,
not only the quality of the final generated response.

\subsubsection{General Capability Retention}
\label{sec:general_capability}
\begin{figure}[!t]
    \centering
    \includegraphics[width=\linewidth]{figures/case_study.pdf}
    \caption{Case study under the same question. \textsc{RetroGen}
produces a more evidence-grounded, less repetitive, and higher-density answer.}
    \label{fig:case_study}
    \vspace{-3mm}
\end{figure}

A potential concern with training on long agentic trajectories is that the model
may become overly specialized to evidence-seeking workflows and lose general
instruction-following or reasoning abilities. Figure~\ref{fig:general_cap} compares each
initial checkpoint with its \textsc{RetroGen}-trained counterpart on GSM8K, MMLU-Pro, IFEval, and HumanEval.
Across backbones, \textsc{RetroGen} shows no systematic degradation: most points are comparable to or slightly above the initial checkpoint. Mixed with a modest amount of general SFT data, verified agentic trajectories therefore improve grounded long-form behavior without inducing catastrophic forgetting.

\subsubsection{Case Study}
\label{sec:case_study}

Figure~\ref{fig:case_study} compares the initial model and \textsc{RetroGen} on the
same question, \emph{``UK Parliament: How does it all work and what are the
numbers?''}. The initial model generates a
plausible overview, but does not attach source links and quickly degenerates
into repeated statements, resulting in low information density.

In contrast, \textsc{RetroGen} produces a more compact and grounded synthesis.
It links claims to the retrieved passages, covers multiple aspects of the
question---including Parliament's role, its two houses, public representation,
and research-supported policymaking---and avoids repetitive phrasing. This
case study illustrates that retrospective process supervision improves not only
evidence use, but also the organization and density of the final answer.
